\famichapter{Problem and Objectives}{This chapter defines the problem that FamiPet
addresses, the objectives of the project, and the deliverables that define a
successful implementation.}{ch:problem}

\section{Problem Statement}

Pet ownership involves several practical responsibilities that are today spread
across separate tools and sources: remembering vaccination schedules, tracking
medical history, keeping appointments with veterinarians, organising daily care
like feeding and medicine, proving a pet's identity when it is lost, and managing
the process of adopting an animal responsibly. Owners keep this information
informally or in unrelated documents, so records are lost, reminders are missed,
and a pet's health history is unavailable at the moment it is needed. People
looking for a pet have similar difficulties: there is no single place to see
animals available for adoption and to apply for them, or to report and search
lost or found animals.

FamiPet addresses this by providing one integrated platform in which:

\begin{itemize}
    \item a user registers, verifies the account, and manages a personal profile;
    \item a user records pets, generating a unique digital pet identity and a QR
    code that encodes public pet information;
    \item health records and vaccinations are kept per pet, with ownership
    enforced by the backend;
    \item veterinary appointments can be booked against a veterinarian's
    availability, with automatic notifications and reminders;
    \item reminders for feeding, medicine, grooming, and care can be created and
    completed;
    \item users can post to a community, like and comment on posts, and report
    lost or found animals;
    \item users can apply to adopt pets, and administrators can approve or reject
    those applications, automatically changing the pet's status; and
    \item an AI assistant answers pet-care questions and gives advice based on
    the user's actual pet records.
\end{itemize}

The problem is therefore the need for an authenticated, backend-driven pet-care
system in which all user data is owned by the correct user, is persisted in a
real database, and is protected by server-side access control.

\section{Objectives}

The objectives of the FamiPet project are:

\begin{enumerate}
    \item To build a full-stack web application for pet care and adoption with an
    authenticating backend and a database-backed data store.
    \item To provide a secure authentication mechanism based on e-mail
    verification, password hashing, and signed bearer tokens.
    \item To enforce ownership and role-based authorization on the backend so
    that users can only access their own resources and only administrators can
    perform administrative actions.
    \item To implement a digital pet identity with a unique pet identifier and a
    QR code that can be read to obtain public pet information.
    \item To implement pet health management, including health records and
    vaccinations, tied to the authenticated owner's pets.
    \item To implement appointment booking with veterinarian slot conflict
    detection and automatic notification and reminder generation.
    \item To implement a reminder system with per-owner, per-pet reminders that
    persist and can be completed.
    \item To implement adoption management with pending, approved, and rejected
    states and automated pet status transitions and notifications.
    \item To implement community and lost-and-found features with moderation by
    administrators.
    \item To integrate an AI pet-care assistant that uses the owner's real pet
    data and falls back to safe, generic guidance when the AI service is
    unavailable.
    \item To validate the system through real end-to-end tests, database
    persistence checks, and security boundary tests.
\end{enumerate}

\section{Deliverables and Success Criteria}

The project deliverables are:

\begin{itemize}
    \item the backend API (Express, MongoDB, Mongoose) at \texttt{backend/};
    \item the frontend application at \texttt{frontend/} and its in-progress
    React migration at \texttt{frontend-react/};
    \item project documentation, including this black book, the implementation
    roadmap (\texttt{ROADMAP.md}, \texttt{docs/migration.md}), and the
    maintenance rules in \texttt{AGENTS.md};
    \item a seeded database containing default breeds, an administrator account,
    and a demo user for evaluation.
\end{itemize}

A phase is considered successful when the functionality is implemented, connects
to the persisted database, passes the associated tests, and is committed and
pushed to the repository. Feature phases are tracked in a checklist: at the time
of writing, the vanilla frontend and backend are feature-complete, and React
migration phases 1--21 are complete while phases 22--29 remain planned
(Chapter~\ref{ch:process}).